\chapter{Teoría local de las curvas en el espacio}

Vamos a trabajar ahora con curvas $\alpha: I \to \bb{R}^3$ p.p.a. con
\begin{gather*}
    \alpha'(s) = T_\alpha(s),\ \ \ |\alpha'(s)| = 1\ \ \ \ \forall s \in I
\end{gather*}
En el caso de curvas planas seguíamos la siguiente secuencia:
\begin{gather*}
    \xymatrix{
        \alpha \ar[r] & T_\alpha \ar[r] & N_\alpha \ar[r] & k_\alpha
    }
\end{gather*}
Sin embargo, para curvas en el espacio no podemos hacer el paso $T_\alpha \to N_\alpha$. Seguiremos entonces la siguiente secuencia (iremos viendo qué son cada uno de estos conceptos):
\begin{gather*}
    \xymatrix{
        \alpha \ar[r] & T_\alpha \ar[r] & k_\alpha \ar[r] & N_\alpha \ar[r] & B_\alpha \ar[r] & \tau_\alpha
    }
\end{gather*}
Definimos así la curvatura de $\alpha$ en un punto $s\in I$ a partir de $T_\alpha$ como
\begin{gather*}
    k_\alpha(s) = |T_\alpha'(s)| \ \ \ \ \forall s \in I
\end{gather*} 
y se verifica que $k_\alpha(s) \geq 0$ para todo $s\in I$, por lo que por definición no habrá curvas en el espacio con curvatura negativa.\\

Defino $N_\alpha(s)$ el vector normal a $\alpha$ en el instante $s$ como 
\begin{gather*}
    N_\alpha(s) = \frac{T'_\alpha(s)}{k_\alpha(s)} = \frac{T'_\alpha(s)}{|T'_\alpha(s)|}
\end{gather*}
A partir de ahora supondremos $k_\alpha(s)>0$ para todo $s\in I$ por lo que siempre estará definido el vector normal 
\begin{gather*}
    s\in I \mapsto \{T_\alpha(s), N_\alpha(s)\}
\end{gather*}
Esto ya no es una base de $\bb{R}^3$ (como ocurría en las curvas planas). Estudiemos las características del espacio que define esta base. Por un lado tenemos que 
\begin{gather*}
    |T_\alpha(s)| = 1 = |N_\alpha(s)|\ \ \ \ \forall s \in I
\end{gather*}
por lo que esta base es normal. Por otro lado tenemos que 
\begin{gather*}
    \langle T_\alpha(s), T'_\alpha(s)\rangle = 0
\end{gather*}
lo que nos dice que esta base es ortogonal. Por tanto $\{T_\alpha(s), N_\alpha(s)\}$ es una base ortonormal de un plano. \\

Sea $B_\alpha:I\to \bb{R}^3$ la función vectorial dada por 
\begin{gather*}
    B_\alpha(s) = T_\alpha(s) \wedge N_\alpha(s)
\end{gather*} 
A esta función la llamaremos el \textbf{vector binormal} de $\alpha$ en el instante $s\in I$.

Hemos construido para todo $s\in I$ un \textbf{triedro de Frenet}
\begin{gather*}
    s \mapsto \{T_\alpha(s), N_\alpha(s), \beta_\alpha(s)\}
\end{gather*}

Buscamos ahora construir las {ecuaciones de Frenet}. Obtenemos los siguientes resultados:
\begin{gather*}
    T'_\alpha(s) = k_\alpha N_\alpha(s)\\
    B'_\alpha = (T_\alpha(s) \wedge N_\alpha(s))' = \cancel{T'_\alpha(s)\wedge N_\alpha(s)} + T_\alpha(s)\wedge N'_\alpha(s) = T_\alpha(s)\wedge N'_\alpha(s)
\end{gather*}
Si desarrollamos ahora 
\begin{gather*}
    |B_\alpha(s)|^2 = 1 \ \ \ \ \forall s \in I \Rightarrow \langle B_\alpha'(s), B_\alpha(s) \rangle = 0 \ \ \ \ \forall s \in I
\end{gather*}
por lo que $B'_\alpha(s)$ no tiene componentes en $B_\alpha(s)$. Como $\langle T_\alpha(s) \wedge N'_\alpha(s), T_\alpha(s)\rangle = 0$ podemos concluir que
\begin{gather*}
    \langle B'_\alpha(s), T_\alpha(s) \rangle = 0
\end{gather*}
Por lo que $B'_\alpha(s)$ tampoco tiene componentes en $T_\alpha(s)$ y podemos concluir que 
\begin{gather*}
    B'_\alpha(s) = \tau_\alpha(s) N_\alpha(s)
\end{gather*}
a dicha $\tau_\alpha$ la llamaremos \textbf{torsión} de $\alpha$ en el instante $s\in I$. Nos faltará ver cómo varía el vector normal. Sabemos que 
\begin{gather*}
    \langle N_\alpha(s), T_\alpha(s) \rangle = 0 \ \ \ \ \forall s \in I
\end{gather*}
Al derivar esta expresión obtenemos 
\begin{gather*}
    \langle N'_\alpha(s), T_\alpha(s) \rangle + \langle N_\alpha(s), T'_\alpha(s) \rangle = 0 \Rightarrow \\
    \Rightarrow \langle N'_\alpha(s), T_\alpha(s) \rangle = - \langle N_\alpha(s), T'_\alpha(s) \rangle = -k_\alpha(s)
\end{gather*}
y ya conocemos la componente de $N'_\alpha$ en $T_\alpha$. Nos acordamos también de que teníamos una base ortonormal por lo que
\begin{gather*}
    \langle N_\alpha(s), B_\alpha(s) \rangle = 0
\end{gather*}
Derivamos esta expresión y tenemos que 
\begin{gather*}
    \langle N'_\alpha (s), B_\alpha(s) \rangle + \langle N_\alpha(s), B'_\alpha(s) \rangle = 0 \Rightarrow \\
    \Rightarrow \langle N'_\alpha (s), B_\alpha(s) \rangle = -\langle N_\alpha(s), B'_\alpha(s) \rangle = -\tau_\alpha(s)
\end{gather*}
y ya tenemos la componente de $N'_\alpha$ en $B_\alpha$. Nos faltaría solo conocer la componente de $N'_\alpha$ en $N_\alpha$. Para ello sabemos que 
\begin{gather*}
    |N_\alpha(s)|^2 = 1
\end{gather*}
y al derivar obtenemos 
\begin{gather*}
    \langle N'_\alpha(s), N_\alpha(s) \rangle = 0
\end{gather*}
y ya tenemos todas sus componentes. Tendremos finalmente:
\begin{gather*}
    N'_\alpha(s) = -k_\alpha(s) T_\alpha(s) - \tau_\alpha(s) B_\alpha(s)
\end{gather*}

De esta forma las \textbf{ecuaciones de Frenet} de una curva p.p.a $\alpha$ serán:
\begin{gather*}
    \left\{
        \begin{array}{l}
            T'_\alpha(s) = k_\alpha N_\alpha(s)\\
            N'_\alpha(s) = -k_\alpha(s) T_\alpha(s) - \tau_\alpha(s) B_\alpha(s)\\
            B'_\alpha(s) = \tau_\alpha(s) N_\alpha(s)
        \end{array}
    \right.
\end{gather*}

\begin{ejercicio}\
    \begin{enumerate}
        \item Supongamos que $\alpha : I \to \bb{R}^3$ es una curva plana p.p.a. Entonces la curvatura de $\alpha$ en cada instante en el espacio coincide con el valor absoluto de la curvatura en cada instante de $\alpha$ considerada como curva plana (considerando $tr(\alpha)\subset P$), es decir, 
        \begin{gather*}
            k^{\bb{R}^3}_\alpha(s) = |k_\alpha^P(s)|
        \end{gather*}

        Comenzamos tomando las ecuaciones de Frenet de $\alpha$ como curva plana:
        \begin{gather*}
            T_\alpha'(s) = k_\alpha^P(s)N_\alpha^P(s)
        \end{gather*}
        Tomamos módulo y nos queda para todo $s\in I$
        \begin{gather*}
            |T'_\alpha (s)| = |k_\alpha^P(s)| \Rightarrow k^{\bb{R}^3}_\alpha(s) = |k_\alpha^P(s)|
        \end{gather*}

        \item Si $\alpha:I\to \bb{R}^3$ es una recta\footnote{o un segmento de recta}, entonces $k_\alpha^{\bb{R}^3} = 0$ (como consecuencia del ejercicio anterior) y no podremos hablar del diedro de Frenet ni de la torsión.
        
        \item Si $\beta : I \to \bb{R}^3$ es una circunferencia\footnote{o un arco de circunferencia} de radio $r$, entonces 
        \begin{gather*}
            k_\beta^{\bb{R}^3}(s) = \frac{1}{r} \ \ \ \ \forall s \in I
        \end{gather*}
        La circunferencia $\beta$ es una curva plana, es decir, que existe un único plano $P\subset \bb{R}^3$ de forma que $tr(\beta)\subset P$. Supongamos que la ecuación del plano $P$ es 
        \begin{gather*}
            \langle x, a \rangle = b ,\ \ \ a\in \bb{R
            }^3\setminus\{0\},\ \ b\in \bb{R}
        \end{gather*}
        Por tanto tendríamos que $\langle \beta(s), a \rangle = b$ para todo $s\in I$. Derivando esta expresión tendremos que para todo $s\in I$
        \begin{gather*}
            \langle T_\beta(s), a\rangle = 0\\
            k_\beta(s)\langle N_\beta(s), a \rangle = 0 \Rightarrow \langle N_\beta(s), a \rangle = 0\\
            \langle N'_\beta(s), a \rangle = 0\\
            \tau_\beta(s) \langle \beta_\alpha(s),a \rangle = 0
        \end{gather*}
        De aquí obtenemos que o bien $B_\beta(s) = a$ o bien $B_\beta(s) = -a$ para todo $s\in I$ lo que (a partir de la tercera ecuación de Frenet) nos dice que 
        \begin{gather*}
            \tau_\beta(s) = 0 \ \ \ \ \forall s \in I
        \end{gather*}

        \item Ahora querremos ver si la condición de $\tau(s)=0$ para todo $s\in I$ se verifica para todas las curvas planas p.p.a. de la forma $\gamma:I \to \bb{R}^3$ tales que $\exists P \subset \bb{R}^3$ con $tr(\gamma)\subset P$ y tales que $k(s)>0$ para todo $k>0$. \\
         
        Sabemos que $T = \gamma'\in \vec{P}$ y por tanto $T'\in \vec{P}$. Como $T$ y $N$ son colineales tendremos que 
        \begin{gather*}
            T,N\in \vec{P} \text{ fijo }
        \end{gather*}
        Lo que nos dice que $B$ es constante y de nuevo tendremos que 
        \begin{gather*}
            B' = \tau N \Rightarrow \tau = 0
        \end{gather*}

        \item[4\ bis.] Consideramos $\alpha :I\to \bb{R}^3$ una curva p.p.a con curvatura $k>0$. Supongamos que $\tau \equiv 0$. Entonces 
        \begin{gather*}
            B'=\tau N \Rightarrow B' \equiv 0 \Rightarrow B \text{ es constante}
        \end{gather*}
        Consideramos $B=v\in \bb{R}^3$ y sabemos que $|v|=1$. Estudiamos ahora la relación entre $T$ y $B$ y llegamos a la conclusión de que son perpendiculares, es decir,
        \begin{gather*}
            \langle T, v \rangle = 0 \Rightarrow \langle \alpha'(s), v \rangle = 0 \ \ \ \forall s \in I
        \end{gather*}
        Notamos que $\langle \alpha'(s), v \rangle = \langle \alpha(s), v \rangle'$ lo que nos dice que
        \begin{gather*}
            \langle \alpha(s), v \rangle = c \ \ \ \ c\in \bb{R}
        \end{gather*}
        Si desarrollamos este producto escalar obtenemos
        \begin{gather*}
            v_1x(s) + v_2y(s) + v_3z(s) = c \ \ \ \ \forall s\in I \text{ donde } (v_1,v_2,v_3)\neq (0,0,0)
        \end{gather*}
        que es la expresión de un plano. Llegamos finalmente a que la traza de $\alpha$ está contenida en un plano.

        \item Veamos ahora el caso de la hélice circular recta. Recordemos que la definíamos como 
        \begin{gather*}
            \alpha(t) = \left(a\cos\left(\frac{t}{\sqrt{a^2+ b^2}}\right), a\sen\left(\frac{t}{\sqrt{a^2+ b^2}}\right), \frac{bt}{\sqrt{a^2 + b^2}}\right)
        \end{gather*}
        con $ab\neq 0$. Sabemos que esta curva no es plana. Comenzamos calculando el vector tangente
        \begin{gather*}
            T(s) = \alpha'(s) = \frac{1}{\sqrt{a^2 + b^2}} \left(-a\sen\left(\frac{t}{\sqrt{a^2+ b^2}}\right), a\cos\left(\frac{t}{\sqrt{a^2+ b^2}}\right), b\right)
        \end{gather*}
        Además, tenemos que 
        \begin{gather*}
            |\alpha'(s)|^2 = \frac{1}{\sqrt{a^2 + b^2}} (a^2 + b^2) = 1 \Rightarrow \text{ p.p.a}
        \end{gather*}
        Si seguimos operando tenemos que 
        \begin{gather*}
            T'(s) = \frac{1}{a^2 + b^2} \left(-a\cos\left(\frac{t}{\sqrt{a^2+ b^2}}\right), -a\sen\left(\frac{t}{\sqrt{a^2+ b^2}}\right), 0\right)
        \end{gather*}
        A partir de esto podemos calcular la curvatura
        \begin{gather*}
            k(s) = |T'(s)| = \frac{|a|}{a^2 + b^2} > 0 \Rightarrow k(s)\text{ es constante}
        \end{gather*}
        Como esta curvatura es positiva podremos calcular el vector normal:
        \begin{gather*}
            N(s) = \frac{T'(s)}{k(s)} = \frac{-a}{|a|} \left(\cos\left(\frac{t}{\sqrt{a^2+ b^2}}\right), \sen\left(\frac{t}{\sqrt{a^2+ b^2}}\right), 0\right)
        \end{gather*}
        Estamos pues en condiciones de calcular el vector binormal.
        \begin{align*}
            B(s) &= T(s) \wedge N(s) = 
            \frac{-a}{|a|(a^2 + b^2)} \begin{vmatrix}
                \vec{i} & \vec{j} & \vec{k}\\
                -a\sen\left(\frac{t}{\sqrt{a^2+ b^2}}\right) & a\cos\left(\frac{t}{\sqrt{a^2+ b^2}}\right) & b\\
                \cos\left(\frac{t}{\sqrt{a^2+ b^2}}\right) & \sen\left(\frac{t}{\sqrt{a^2+ b^2}}\right) & 0
            \end{vmatrix} = \\
            &=\frac{-sgn(a)}{(a^2 + b^2)} \left(-b \sen \left(\frac{t}{\sqrt{a^2+ b^2}}\right), b\cos \left(\frac{t}{\sqrt{a^2+ b^2}}\right), -a\right) 
        \end{align*}
        Derivamos el vector binormal en vistas de calcular la torsión
        \begin{gather*}
            B'(s) = \frac{-sgn(a)b}{a^2 + b^2} \left(- \cos \left(\frac{t}{\sqrt{a^2+ b^2}}\right), -\sen \left(\frac{t}{\sqrt{a^2+ b^2}}\right), 0\right) 
        \end{gather*}
        y llegamos a que $B\| N$, es decir que son colineales y en concreto 
        \begin{gather*}
            \tau_\alpha(s) = \frac{-b}{a^2 + b^2}
        \end{gather*}
        que nos dice que es constante no nula y además nos confirma que la dirección del giro solo depende de $b$.

        \item Sea $\alpha: I \to \bb{R}^3$ p.p.a con $k(s)>0$. Supongamos que $k(s)=k_0$ (constante) y $\tau(s) = \tau_0$ (constante también) para todo $s\in I$. Queremos ver que entonces que $tr(\alpha)$ está contenida en una hélice circular recta.\\
        
        Defino $a,b \in \bb{R}$ tales que 
        \begin{gather*}
            k_0 = \frac{|a|}{a^2 + b^2},\ \ \ \
            \tau_0 = \frac{-b}{a^2 + b^2}
        \end{gather*}
        Tendremos que ver que existen soluciones a estas ecuaciones. Elevando al cuadrado tenemos 
        \begin{gather*}
            k_0^2 = \frac{a^2}{(a^2 + b^2)^2},\ \ \ \
            \tau_0^2 = \frac{b^2}{(a^2 + b^2)^2}
        \end{gather*}
        Uniendo estas expresiones llegamos a que 
        \begin{gather*}
            a^2 + b^2 = \frac{1}{k_0^2 + \tau_0^2}
        \end{gather*}
        Despejando ahora de las ecuaciones iniciales tenemos 
        \begin{gather*}
            |a| = \frac{k_0}{k_0^2 + \tau_0^2},\ \ \ \ b = - \frac{\tau_0}{k_0^2 + \tau_0^2}
        \end{gather*}
        Sea $h:I\to \bb{R}^3$ una hélice circular recta con parámetros $a$ y $b$. Entonces, del ejemplo anterior tenemos que 
        \begin{gather*}
            k_h(s) = \frac{a}{a^2 + b^2} = k_0\ \ \ \ \forall s \in I\\
            t_h(s) = - \frac{b}{a^2 + b^2}\ \ \ \ \forall s \in I
        \end{gather*}
        Ahora $\alpha$ y $h$ son dos curvas p.p.a. con 
        \begin{gather*}
            k_\alpha = k_h = k_0,\ \ \ \ \tau_\alpha = \tau_h = \tau_0
        \end{gather*}
        Sea $s_0\in I$ e impongo las siguientes condiciones:
        \begin{gather*}
            \alpha(s_0)=h(s_0)\\
            T_\alpha(s_0) = T_h(s_0),\ \ \ \ N_\alpha(s_0) = N_h(s_0),\ \ \ \ B_\alpha(s) = B_h(s_0)
        \end{gather*}
        Defino $f:I\to \bb{R}$ dada por 
        \begin{gather*}
            f(s) = \frac{1}{2} \{|T_\alpha(s) - T_h(s)|^2 + |N_\alpha(s) - N_h(s)|^2 + |B_\alpha(s) - B_h(s)|^2\}
        \end{gather*}
        Además, por las condiciones impuestas tenemos que $f(s_0)=0$. Al derivar $f$ obtenemos 
        \begin{align*}
            f' &= \langle T'_\alpha - T'_h, T_\alpha - T_h\rangle + \langle N'_\alpha - N'_h, N_\alpha - N_h\rangle + \langle B'_\alpha - B'_h, B_\alpha - B_h\rangle = \\
            &= k_0 \langle N_\alpha - N_h , T_\alpha - T_h\rangle - k_0\langle T_\alpha - T_h, N_\alpha-N_h\rangle - \tau_0 \langle B_\alpha - B_h, N_\alpha-N_h \rangle + \tau_0\langle N_\alpha - N_h, B_\alpha - B_h \rangle =\\
            &= 0
        \end{align*}
        Obtenemos por tanto que $f(s_0)=0$ y $f'(s)=0$ para todo $s\in I$ lo que nos dice que $f(s)=0$ para todo $s\in I$. Por la construcción de $f$ tenemos que 
        \begin{gather*}
            \alpha' = T_\alpha = T_h = h' \Rightarrow (\alpha - h)'(s) = 0\ \ \ \ \forall s \in I
        \end{gather*}
        Con la condición de $\alpha(s_0)=h(s_0)$ tenemos que $\alpha(s)=h(s)$ para todo $s\in I$.

        \item (Siempre lo pone en los exámenes) Consideramos la curva $\alpha:I\to \bb{R}^3$ p.p.a con $k(s)>0$ y $\tau(s)>0$ para todo $s\in I$. Si $tr(\alpha)$ está contenida en una esfera $S^2(a,R)$ con $a\in \bb{R}^3$ y $R>0$. Nos preguntamos si existe alguna relación entre $k$ y $\tau$.\\
        
        Dado un $s\in I$ tenemos que $\alpha(s)\in S^2(a,R)$ si y solo si $|\alpha(s) - a|^2 = R^2$. Tenemos por tanto
        \begin{gather*}
            \langle T, \alpha - a \rangle = 0
        \end{gather*}
        Derivando obtenemos 
        \begin{gather*}
            k\langle N , \alpha - a\rangle + 1 = 0 \Rightarrow \langle N, \alpha - a \rangle = - \frac{1}{k}\\
            \langle N', \alpha - a \rangle = - \left(\frac{1}{k}\right)' \Rightarrow \tau \langle B, \alpha -a \rangle = \left(\frac{1}{k}\right)'
        \end{gather*}
        Por lo que llegamos a que
        \begin{gather*}
            \langle B, \alpha -a \rangle = \frac{1}{\tau} \left(\frac{1}{k}\right)'
        \end{gather*}
        Tenemos las tres componentes del vector $\alpha -a$ respecto de la base ortonormal $\{T,N,B\}$:
        \begin{gather*}
            \alpha - a = - \frac{1}{k} N + \frac{1}{\tau} \left(\frac{1}{k}\right)'B
        \end{gather*}
        Tomándo módulo 
        \begin{gather*}
            |\alpha - a|^2 = \left(\frac{1}{k}\right)^2 + \frac{1}{\tau^2} \left[\left(\frac{1}{k}\right)'\right]^2
        \end{gather*} 
        Finalmente llegamos a que 
        \begin{gather*}
            R^2 = \left(\frac{1}{k}\right)^2 + \frac{1}{\tau^2} \left[\left(\frac{1}{k}\right)'\right]^2
        \end{gather*}

        \item[7 bis.] Nos preguntamos ahora si esta condición es suficiente para determinar que la traza de una curva está contenida en una circunferencia. Sea $\alpha: I \to \bb{R}^3$ una curva p.p.a con $k(s)>0$, $\tau(s)>0$ y $k'(s)>0$ para todo $s\in I$ (esta última condición es \textit{ad-hoc}). Supongamos que se verifica
        \begin{gather*}
            \left(\frac{1}{k}\right)^2 + \left[\frac{1}{\tau}\left(\frac{1}{k}\right)'\right]^2 = C \in \bb{R}
        \end{gather*}
        Entonces la traza de $\alpha$ está contenida en una esfera.\\

        Si $C>0$ tenemos que existe un $R>0$ tal que $C = R^2$. Podemos traducir la condición en 
        \begin{gather*}
            \left(\frac{1}{k}\right)^2 + \left[\frac{1}{\tau}\left(\frac{1}{k}\right)'\right]^2 = R^2
        \end{gather*}
        Del apartado anterior tenemos que definir una función
        \Func{f}{I}{\bb{R}}{s}{\alpha + \frac{1}{k}N - \frac{1}{\tau}\left(\frac{1}{k}\right)'B}
        Veamos ahora que esta función es constante. Para ello calculamos 
        \begin{align*}
            f' &= T + \left(\frac{1}{k}\right)' N + \frac{1}{k}(-k T - \tau B) - \left(\frac{1}{\tau}\left(\frac{1}{k}\right)'\right)'B - \left(\frac{1}{k}\right)' N =\\
            &= - \left[ \frac{\tau}{k} + \left(\frac{1}{\tau}\left(\frac{1}{k}\right)'\right)' \right]B
        \end{align*}
        Derivamos ahora la condición del enunciado buscando seguir desarrollando la derivada de $f$:
        \begin{gather*}
            2 \left(\frac{1}{k}\right) \left(\frac{1}{k}\right)' + 2 \left[\frac{1}{\tau} \left(\frac{1}{k}\right)'\right]\left[\frac{1}{\tau} \left(\frac{1}{k}\right)'\right]' = 0
        \end{gather*}
        Multiplicando por $\frac{\tau}{2}$ tenemos 
        \begin{gather*}
            \frac{\tau}{k} \left(\frac{1}{k}\right)' +  \left(\frac{1}{k}\right)'\left[\frac{1}{\tau} \left(\frac{1}{k}\right)'\right]' = \\
            = \left(\frac{1}{k}\right)' \left(\frac{\tau}{k} + \left[\frac{1}{\tau} \left(\frac{1}{k}\right)' \right]'\right) = 0
        \end{gather*}
        Como $(\frac{1}{k})' = -\frac{k'}{k^2} < 0$ por la condición \textit{ad-hoc} del enunciado nos queda que necesariamente
        \begin{gather*}
            \frac{\tau}{k} + \left[\frac{1}{\tau} \left(\frac{1}{k}\right)' \right]' = 0
        \end{gather*}
        Lo que finalmente, al sustituir en $f'$ nos dice $f'\equiv 0$, o equivalentemente, $f(s) = a$ para todo $s\in I$, es decir
        \begin{align*}
            a &= \alpha + \frac{1}{k}N - \frac{1}{\tau}\left(\frac{1}{k}\right)'B\\
            a-\alpha &= \frac{1}{k}N - \frac{1}{\tau}\left(\frac{1}{k}\right)'B\\
            |\alpha - a|^2 &= \frac{1}{k^2} - \left[\frac{1}{\tau}\left(\frac{1}{k}\right)'\right]^2 = R^2
        \end{align*}
        Lo que nos dice que $tr(\alpha)\subset S^2(a, R)$.

        \item Sea $\alpha : I \to \bb{R}^3$ una curva p.p.a con $k(s)>0$. Entonces $\alpha$ es un arco de circunferencia si y solo si tiene $k$ constante y $tr(\alpha)$ contenida en una esfera.
        
        \item Sea $\alpha:I\to \bb{R}^3$ una curva p.p.a con $k(s)>0$. Si $tr(\alpha)$ está contenida en una esfera y tiene torsión constante $a\in \bb{R}$, probar que existen $b,c\in \bb{R}$ tales que 
        \begin{gather*}
            k(s) = \frac{1}{b\cos(as) + c\cos(as)}
        \end{gather*}

        \item Sea $ \alpha : I \to \bb{R}^3$ una curva regular. Consideramos 
        \begin{gather*}
            \xymatrix{
                I \ar@<1ex>[r]^{S_a} & J \ar@<1ex>[l]^{h=S_a^{-1}}
            }\ \ \ \ \beta = \alpha \circ h \text{ reparametrización de }\alpha
        \end{gather*}
        Definimos la curvatura de $\alpha$ por 
        \begin{gather*}
            k_\alpha(t) = k_\beta(S_a(t))
        \end{gather*}
        Y si $k_\alpha(t)>0$ para todo $t\in I$, definimos su torsión por 
        \begin{gather*}
            \tau_\alpha(t) = \tau_\beta(S_a(t))
        \end{gather*}
        Demostrar que 
        \begin{gather*}
            k_\alpha(t) = \frac{|\alpha'(t) \wedge \alpha''(t)|}{|\alpha'(t)|^3}\ \ \ \ \ \ \ \tau_\alpha(t) = - \frac{\det(\alpha'(t), \alpha''(t), \alpha'''(t))}{|\alpha'(t) \wedge \alpha''(t)|^2}
        \end{gather*}
    \end{enumerate}
\end{ejercicio}

\section{Teorema Fundamental de la Teoría Local de Curvas en el Espacio}
\begin{teo}[Teorema Fundamental de la Teoría Local de Curvas en el Espacio]
    Sean $k_0, \tau_0:I\to \bb{R}$ las funciones derivables con $k_0>0$. Entonces existe una curva $\alpha:I\to \bb{R}^3$ p.p.a. y tal que 
    \begin{gather*}
        k_\alpha(s) = k_0(s)\ \ \ \ \tau_\alpha(s)=\tau_0(s)\ \ \ \ \forall s \in I
    \end{gather*}
    Además, $\alpha$ es única salvo movimiento rígido directo
    \begin{proof}
        Defino $A_0:I\to M_9(\bb{R})$ como 
        \begin{gather*}
            A_0 = \left(\begin{array}{c|c|c}
                0_3 & k_0I_3 & 0_3\\
                \hline
                -k_0I_3 & 0_3 & -\tau_0I_3\\
                \hline
                0_3 & \tau_0I_3 & 0_3
            \end{array}\right)
        \end{gather*}
        Como esa función $A_0$ planteo la siguiente EDO\footnote{Ecuación Diferencial Ordinaria (no es ordinaria porque diga palabrotas, es porque solo tiene una derivada).}
        \begin{gather*}
            x'=A_0x
        \end{gather*}
        donde $x:I\to \bb{R}^9$. Sabemos que si $a\in \bb{R}^9$, entonces existe una única solución $f$ de la EDO tal que $f(s_0)=a$ para cierto $s_0\in I$ y $f$ está definida en todo $I$. \\

        Elijo $a=(a_1, a_2, a_3, a_4, a_5, a_6, a_7, a_8, a_9)$ de forma que 
        \begin{gather*}
            t_0 = \begin{pmatrix}
                a_1 \\
                a_2\\
                a_3
            \end{pmatrix} \ \ \ \ \ \ 
            n_0 = \begin{pmatrix}
                a_4 \\
                a_5\\
                a_6
            \end{pmatrix} \ \ \ \ \ \ 
            b_0 = \begin{pmatrix}
                a_7 \\
                a_8\\
                a_9
            \end{pmatrix}
        \end{gather*}
        sea una base ortonormal positivamente orientada de $\bb{R}^3$. Para cada $s\in I$ defino 
        \begin{gather*}
            t(s) = \begin{pmatrix}
                f_1 \\
                f_2\\
                f_3
            \end{pmatrix}(s) \ \ \ \ \ \ 
            n(s) = \begin{pmatrix}
                f_4 \\
                f_5\\
                f_6
            \end{pmatrix}(s) \ \ \ \ \ \ 
            b(s) = \begin{pmatrix}
                f_7 \\
                f_8\\
                f_9
            \end{pmatrix}(s)
        \end{gather*}
        Entonces 
        \begin{gather*}
            \begin{pmatrix}
                t'\\
                n'\\
                b'
            \end{pmatrix} = f' = A_0f = A_0 \begin{pmatrix}
                t\\
                n\\
                b
            \end{pmatrix}
        \end{gather*}
        Nos quedarán las siguientes ecuaciones junto con las condiciones iniciales
        \begin{gather*}
            \left\{
                \begin{array}{l}
                    t'=k_0n\\
                    n'=-k_0t - \tau_0 b \\
                    b' = \tau n
                \end{array}
            \right. \ \ \ \ \ \ 
            \left\{
                \begin{array}{l}
                    t(s_0) = t_0\\
                    n(s_0) = n_0\\
                    b(s_0) = b_0
                \end{array}
            \right.
        \end{gather*}
        Defino $M:I \to M_3(\bb{R})$ por 
        \begin{gather*}
            M = \begin{pmatrix}
                |t|^2 & \langle t,n \rangle & \langle t,b \rangle \\
                \langle n,t \rangle & |n|^2 & \langle n,b \rangle \\
                \langle b,t \rangle & \langle b,n \rangle & |b|^2
            \end{pmatrix}
        \end{gather*}
        Se deja como ejercicio demostrar que 
        \begin{gather*}
            M' = AM - MA
        \end{gather*}
        que es de nuevo una EDO distinta donde $A$ está definida por 
        \begin{gather*}
            A = \begin{pmatrix}
                0 & k & 0\\
                -k_0 & 0 & -\tau_0\\
                0 & \tau_0 & 0
            \end{pmatrix}
        \end{gather*}
        Tenemos entonces que $M$ es solución de la segunda EDO con condición inicial
        \begin{gather*}
            M(s_0) = I_3
        \end{gather*}
        pero hay otra solución de la segunda EDO con la misma condición inicial, $M_0(s) = I_3$ para todo $s$. Por la unicidad de soluciones de EDOs para misma ecuación y condición inicial necesariamente tendremos que 
        \begin{gather*}
            M(s) = I_3 \ \ \ \forall s \in I
        \end{gather*}
        Entonces $t(s)$, $n(s)$, $b(s)$ es una base ortonormal de $\bb{R}^3$ para todo $s\in I$ por lo que $det(t(s), n(s), b(s)) = \pm 1$. Como $det(t(s_0), n(s_0), b(s_0)) = 1$ y el determinante no puede saltar, necesariamente $t(s)$, $n(s)$ y $b(s)$ forman una base ortonormal positivamente orientada de $\bb{R}^3$ para todo $s\in I$.
    \end{proof}
\end{teo}